\section*{Chapter \ref{ch:b2:corrosion} --- Corrosion and Protection}

\begin{solution}{exo:b2:corrosion:1}
$j = \qty{0.020}{A/m^2}$: $de/dt = 0.020 \times 65.38/(2 \times 96\,485 \times 7.13
\times 10^6) = \qty{9.5e-13}{m/s}$, that is \qty{0.030}{mm} per year.
\end{solution}

\begin{solution}{exo:b2:corrosion:2}
Steel bolts on copper: iron, the less noble, is the anode, and the small bolts
on a large copper cathode corrode fast. Zinc nails in steel: the zinc corrodes
and protects the steel around it.
\end{solution}

\begin{solution}{exo:b2:corrosion:3}
In the narrow gap between the plates and under the bolt head, where the water
is renewed slowly and its oxygen is soon used up: by \omterm{def:b2:corrosion:differential}{differential aeration}
these confined zones become anodes, the open surfaces cathodes.
\end{solution}

\begin{solution}{exo:b2:corrosion:4}
The tin can: tin, nobler than iron, makes the exposed steel the anode of a
large cathode. On the galvanised bucket the zinc corrodes instead of the
steel.
\end{solution}

\begin{solution}{exo:b2:corrosion:5}
Stirring thins the \omterm{def:b2:current-potential-curves:limiting-current}{diffusion layer}: the oxygen plateau rises, the anodic curve
of iron meets it higher, so both the \omterm{def:b2:corrosion:uniform}{corrosion potential} and the \omterm{def:b2:corrosion:uniform}{corrosion
current} increase. Without oxygen, the only cathodic reaction left in neutral
water is the slow reduction of water: the \omterm{def:b2:corrosion:uniform}{corrosion potential} falls and the
\omterm{def:b2:corrosion:uniform}{corrosion current} becomes very small.
\end{solution}

\begin{solution}{exo:b2:corrosion:6}
$t = 0.90 \times 5000 \times 2 \times 96\,485/(65.38 \times 0.10) = \qty{1.33e8}{s}$,
about \qty{4.2}{years}.
\end{solution}

\begin{solution}{exo:b2:corrosion:7}
Charge per kilogram, $nF/M$: zinc \qty{820}{A.h/kg}, magnesium \qty{2.21e3}{A.h/kg},
aluminium \qty{2.98e3}{A.h/kg}. Standard potentials: \ce{Zn^2+}/\ce{Zn} $-0.76$,
\ce{Mg^2+}/\ce{Mg} $-2.36$, \ce{Al^3+}/\ce{Al} \qty{-1.68}{V}. In a dry, resistive soil
the current must be driven through a large resistance: magnesium, much further
below iron, offers the largest driving voltage.
\end{solution}

\begin{solution}{exo:b2:corrosion:8}
$U = 2.0 \times 4.0 + 1.5 = \qty{9.5}{V}$; power \qty{19}{W}; per year $19 \times
8766 = \qty{1.7e5}{W.h}$, \qty{1.7e2}{kW.h}.
\end{solution}

\begin{solution}{exo:b2:corrosion:9}
At the water line the metal is wetted by water rich in oxygen, a cathode; just
below, the water is less aerated and the zone becomes the anode of the cell
formed with the water line: the steel corrodes there, in a band.
\end{solution}

\begin{solution}{exo:b2:corrosion:10}
In aerated water the \omterm{def:b2:corrosion:uniform}{corrosion potential} of stainless steel falls on its
passive plateau: a tiny current, the film healing itself. Chloride ions break
the film at weak points; the bare metal there is a small anode facing an
enormous passive cathode, so its current density is huge. Inside the pit the
metal ions hydrolyse and acidify the solution, chloride is drawn in to balance
the charge, and the film cannot reform: the pit deepens.
\end{solution}

\begin{solution}{exo:b2:corrosion:11}
Oxygen is reduced on all \qty{22}{cm^2}: \omterm{def:b2:current-potential-curves:ie-curve}{cathodic current} $22 \times 5 =
\qty{110}{\micro A}$, all supplied by the oxidation of the \qty{2}{cm^2} of steel
near the joint, $j = \qty{55}{\micro A/cm^2}$. The loss is $5.5$ times that of iron
at \qty{10}{\micro A/cm^2}, about \qty{0.64}{mm} per year.
\end{solution}

\begin{solution}{exo:b2:corrosion:12}
Iron is immune below the \ce{Fe^2+}/\ce{Fe} line, $-0.41 + 0.0296\log 10^{-6} =
\qty{-0.59}{V}$: held below about \qty{-0.6}{V} the steel no longer dissolves. Much
lower, water is reduced at a growing rate on the steel (the \ce{H+}/\ce{H2} line is at
\qty{-0.41}{V} at pH 7 and the \omterm{def:b2:current-potential-curves:overpotential}{overpotential} on steel is a few tenths of a volt):
current is wasted and hydrogen lifts the coating.
\end{solution}

\begin{solution}{pb:b2:corrosion:1}
\textbf{1.} \ce{Fe -> Fe^2+ + 2e-}; \ce{O2 + 2H2O + 4e- -> 4OH-}.
\textbf{2.} The reduction of oxygen, slow and reaching its diffusion plateau,
limits the \omterm{def:b2:current-potential-curves:ie-curve}{cathodic current}; the anodic curve of iron meets it on that plateau.
\textbf{3.} $j = \qty{0.10}{A/m^2}$: $0.10 \times 3.156 \times 10^7/(2 \times 96\,485)
\times 55.845 = \qty{913}{g}$ per square metre and per year.
\textbf{4.} $913/7.87 \times 10^6 = \qty{1.16e-4}{m}$, \qty{0.12}{mm} per year.
\textbf{5.} $4/0.116 = \qty{34}{years}$.
\textbf{6.} The corrosion is not uniform: crevices, differences of aeration
between soils, and defects of the coating concentrate the \omterm{def:b2:current-potential-curves:ie-curve}{anodic current} on
small areas, where pits perforate the wall long before.
\textbf{7.} $E^\circ = \Delta_f G^\circ/(nF)$ for each couple: \ce{Fe} $-0.41$, \ce{Zn} $-0.76$,
\ce{Mg} $-2.36$, \ce{Al} \qty{-1.68}{V}.
\textbf{8.} Zinc, magnesium and aluminium, all below iron.
\textbf{9.} Aluminium passivates: its oxide film stops its dissolution and it
then delivers little current, unless alloyed to prevent the film (which works
in sea water, not in soils).
\textbf{10.} The current must cross the resistance of the soil: the driving
voltage between the anode and the protected steel is about \qty{1.9}{V} for
magnesium against \qty{0.35}{V} for zinc (from the standard potentials).
\textbf{11.} $3.156 \times 10^7 \times 24.305/(2 \times 96\,485 \times 0.50) =
\qty{7.95}{kg}$ per ampere and per year.
\textbf{12.} $\pi \times 0.30 \times 10\,000 = \qty{9.42e3}{m^2}$.
\textbf{13.} \qty{94.2}{m^2}.
\textbf{14.} $0.020 \times 94.2 = \qty{1.88}{A}$.
\textbf{15.} $1.885 \times 20 \times 3.156 \times 10^7 = \qty{1.19e9}{C}$.
\textbf{16.} $1.19 \times 10^9 \times 24.305/(2 \times 96\,485 \times 0.50) =
\qty{3.00e5}{g}$, \qty{300}{kg} (or $1.885 \times 20 \times 7.95$).
\textbf{17.} $300/15 = 20$ anodes.
\textbf{18.} The resistance of the soil limits how far a single anode can push
its current along the pipe: distributed anodes protect it evenly.
\textbf{19.} A rectifier drives current from an inert anode buried in a ground
bed, through the soil, into the pipe connected to its negative pole.
\textbf{20.} $1.885 \times 5.0 + 1.5 = \qty{10.9}{V}$.
\textbf{21.} $10.9 \times 1.885 = \qty{20.6}{W}$; $20.6 \times 8766 = \qty{1.8e5}{W.h}$,
\qty{180}{kW.h} per year.
\textbf{22.} Water is reduced at the pipe: hydrogen forms, lifts the coating and
can embrittle the steel, and current is wasted.
\textbf{23.} $E < -0.41 + 0.0296\log(10^{-6}) = \qty{-0.59}{V}$.
\textbf{24.} Twenty years of protection need $\boldsymbol{\approx \qty{3.0e2}{kg}}$
of magnesium anodes.
\end{solution}
